\chapter{Glossário de Termos Técnicos}
\label{apendice:glossario}

Este apêndice apresenta uma lista de definições para os principais termos técnicos utilizados ao longo deste trabalho.

\begin{description}
    \item[Análise Multicanal] Abordagem de modelagem que decompõe um processo cíclico (como o consumo diário de energia) em múltiplos canais, onde cada canal representa um período de tempo específico (e.g., uma hora do dia). A análise estatística é então realizada intra-canais, comparando medições de um determinado período apenas com dados históricos do mesmo período, aumentando a precisão na detecção de anomalias~\cite{Braga2010TeseValidacao}.

    \item[Building Management System (BMS)] Um sistema de controle centralizado que monitora e gerencia os equipamentos mecânicos e elétricos de um edifício, como ventilação, iluminação e sistemas de energia~\cite{Lee2013BMS}.

    \item[Controle Estatístico de Processo (CEP)] Metodologia estatística utilizada para monitorar, controlar e otimizar processos por meio da análise de sua variabilidade. O CEP ajuda a distinguir entre variações comuns (inerentes ao processo) e causas especiais (eventos não aleatórios) que podem indicar uma mudança no processo~\cite{Braga2013SPCBuildingEnergy}.

    \item[Edge Computing] Um paradigma de computação distribuída que aproxima a computação e o armazenamento de dados das fontes de dados. Em vez de enviar dados para uma nuvem centralizada para processamento, o processamento ocorre localmente, na "borda" da rede. Isso reduz a latência e o consumo de largura de banda, sendo benéfico para aplicações que exigem respostas em tempo real~\cite{Shi2016Edge}.

\end{description}
